
résolu de
façon satisfaisante la question soulevée au n\textsuperscript{o} 1. Comme support pour la formule (4.5) dans $A(X)$ lui-même, signalons le fait qu'il est clair à priori que la restriction du premier membre à $X-Y$ est nulle, donc, $Y$ étant supposé irréductible, le premier membre doit être à priori un multiple entier de $\mathrm{c}\ell(Y)$:
\[
 c_i(Y)=n(Y)\mathrm{c}\ell(Y).
\]
Cela prouve déjà la formule exacte (4.5) lorsque $\mathrm{c}\ell(Y)$ n'est pas un élément de torsion de $A^i(X)$. Il est immédiat que tel est le cas si $X$ est projectif (et non seulement quasi-projectif) sur le corps $k$, en notant que si $\xi\in A^1(X)$ est une polarisation de $X$, alors $\deg(Y\cdot\xi^{\dim Y})>0$; cela prouve la formule voulue dans ce cas. Il en résulte que cette formule est encore valable chaque fois que $X$ est isomorphe à un sous-schéma ouvert d'un schéma projectif lisse sur $k$, ce qui sera le cas lorsqu'on dispose d'une solution convenable du problème de résolution des singularités. Il en est ainsi en particulier si $k$ est de caractéristique nulle, grâce à Hironaka [11].

La validité de la formule exacte (4.5) donc de (4.4) impliquerait que le noyau de $\phi^i:A^i(X)\longrightarrow \Gtop^i(X)$ est annulé par $(i-1)!$, et en particulier \uline{que} $\phi^2:A^2(X)\to\Gtop^2(X)$ \uline{est un isomorphisme}. (Que $\phi^1$ soit un isomorphisme résulte déjà immédiatement de X 5.3.2). Sauf erreur, le fait que $\phi^2$ soit un isomorphisme n'est pas prouvé à l'heure actuelle, même dans le cas particulier où $\dim X=3$. C'est hâtivement en tous cas que le rédacteur avait affirmé la validité de (4.4) et (4.5) dans [9, n\textsuperscript{o}~4, 3\textsuperscript{o}].

\textbf{4.4.} A l'heure actuelle, le principal avantage de $\Gtop^\bullet(X)$ (ou de $\operatorname{Gr}^\bullet(X)$) sur l'anneau de Chow $A(X)$ est le fait qu'il se prête mieux aux calculs explicites de certaines formules d'intersection entre cycles, dans le cas d'intersections ``excédentaires'' i.e. de dimension trop grande. La définition de la classe d'intersections, à l'aide du lemme de Chow permettant de bouger suffisamment un cycle dans sa classe, se prête en effet malaisément au calcul. Un exemple frappant à cet égard est la formule fondamentale suivante [Exp. VII 2.7], pour une immersion fermée $f:Y\to X$ de schémas quasi-projectifs lisses sur $k$:
\begin{equation}\tag{4.6}
 f^*f_*(y)=y\,c_d(\underline{N}),
\end{equation}
formule valable dans $\Gtop^\bullet(Y)$, $y$ étant dans $\Gtop^\bullet(Y)$, $\underline{N}$ le faisceau conormal 
